\documentclass[11pt,letterpaper]{article}
\input{headings}
\newcommand \recipeName {Chicken Filling}
\newcommand \fileName {ChickenFilling}
\chead{\recipeName}

\begin{document}
\input{title}

This is a general-purpouse beef filling recipe that could be used for Pastel\~ao, risolis, or past\'eis. It is also delicious in a simple pressed sandwich or as a topping for pizza.
 
\begin{description}

\item[Ingredients:]\ \\
	\begin{itemize}
	\item 1/2 pound of regular or lean ground beef
	\item 1 medium onion (about 1 cup when finely diced)
	\item 1 cup of drained can whole tomatoes chopped
	\item 1 1/2 Tablespoon of Better Than Bullion concentrated chicken stock
        \item Italian parsley
	\end{itemize}

\item[Procedure:]\ \\
	\begin{enumerate}
	\item[Make the Filling]
	\begin{itemize}
        \item Cook the ground beef until all the water evaporates and it starts sizzling. Do not sautee for browning.
        \item Place a strainer over a large bowl and put the ground beef with its melted fat in the strainer. Let it sit for a few minutes for the fat to drain out of the ground beef.
        \item Return two Tablespoons of the beef rendered fat to the pan where you cooked the ground beef.
        \item Sautee the onions for a few minutes until it is translucent and just very lightly golden. 
        \item Add tree Tablespoons of water to the onions, cover the pan and cook for a few minutes to steam the onions. 
        \item Remove the lead and continue cooking until all the water has evaporated.
        \item Add the chopped tomatoes and the concentrated chicken stock.
        \item Cook for between 5 and 10 minutes until the tomatoes soften.
        \item Return the drained ground beef to the pot.
        \item Either discard or save the remaining rendered fat for a different application.
        \item Cook for about 10 minutes, stirring from time to time.
        \item Taste for salt and adjust to your taste. The filling should be slightly salty.
        \item The filling should be dry with the onions and tomatoes incorporated into the ground beef.
        \end{itemize}
  	\end{enumerate}
\end{description}

\input{\imageDir/\fileName/imageTable}

\end{document}



